\subsection{Bounded subsets in certain inductive limits}

PROPOSITION 5. --- Let \((\mathrm{E}_{i})_{i\in\mathbf{I}}\) be a family of Hausdorff locally convex spaces, and let E be the topological direct sum of this family (II, p.~29). In order that a subset B of E be bounded, it is necessary and sufficient that there exists a finite subset J of I such that \(\operatorname{pr}_{i}(\mathbf{B})\) is bounded in \(E_{i}\) for \(i\in J\) and \(\operatorname{pr}_{i}(\mathbf{B})\subset\{0\}\) for all \(i\notin J\) .

Let J be a finite subset of I. Since the topology of E induces the product topology on \(\prod_{j\in J}E_{j}\) (II, p.~30, prop. 7 and p.~31, prop. 8), it follows from III, p.~4, cor. 2 that the condition is sufficient.

Conversely, let B be a bounded subset of E. Then \(\mathrm{pr}_{i}(\mathbf{B})\) is bounded for all i (III, p.~4, cor. 1). Therefore it is enough to prove that there exists a finite subset J of I such that \(\mathrm{pr}_{i}(\mathbf{B}) \subset \{0\}\) for all \(i \notin J\) . If not, then there exists an infinite sequence \((i_{n})\) of distinct elements of I and an infinite sequence \((x_{n})\) of elements of B such that \(\mathrm{pr}_{i_{n}}(x_{n}) \neq 0\) . Since \(E_{i_{n}}\) is Hausdorff, there exists a continuous semi-norm \(p_{n}\) on \(E_{i_{n}}\) such that \(p_{n}(\mathrm{pr}_{i_{n}}(x_{n})) \geqslant n\) . Hence \(p = \sum_{n \geqslant 1} p_{n} \circ \mathrm{pr}_{i_{n}}\) is a continuous semi-norm on E and p is not bounded on B, which is a contradiction.

PROPOSITION 6. --- Let E be a locally convex space which is the strict inductive limit of an increasing sequence \((E_{n})\) of closed vector subspaces of E (II, p.~33). A subset B of E is bounded if and only if it is contained in one of the subspaces \(E_{n}\) , and is bounded in this subspace.

The condition is sufficient, since the topology induced on \(E_{n}\) by that of E is precisely the given topology of \(E_{n}\) (II, p.~32, prop. 9). To see that the condition is necessary, it is enough (III, p.~4, prop. 3) to prove that if a sequence \((x_{m})\) of points of E is not contained in any of the subspaces \(E_{n}\) , then it cannot tend to 0. By extracting a subsequence of the sequence \((x_{m})\) , we can assume that there exists a strictly increasing sequence \((n_{k})\) of integers such that, for every index k, we have \(x_{k} \notin E_{n_{k}}\) and \(x_{k} \in E_{n_{k+1}}\) . Then there exists (II, p.~33, lemma 2) an increasing sequence \((V_{k})\) of convex sets such that \(V_{k}\) is a neighbourhood of 0 in \(E_{n_{k}}\) , \(V_{k+1} \cap E_{n_{k}} = V_{k}\) and such that \(x_{k} \notin V_{k+1}\) for every index k. The union V of the \(V_{k}\) is then a neighbourhood of 0 in E, and we have \(x_{k} \notin V\) for all k. This proves that the sequence \((x_{k})\) does not tend to 0.

The conclusion of prop. 6 is not necessarily true for a space \(\mathbf{E}\) which is the inductive limit of a non-denumerable directed set of closed subspaces of \(\mathrm{E}\) (cf.~III, p.~38, exerc. 7).

PROPOSITION 7. --- Let \((\mathrm{E}_{n})_{n\geq0}\) be a sequence of Hausdorff locally convex spaces, and for every n, let \(u_{n}:E_{n}\to E_{n+1}\) be an injective linear mapping which is compact (i.e.~such that there exists a neighbourhood of 0 in \(E_{n}\) whose image under \(u_{n}\) is relatively compact; this implies that \(u_{n}\) is continuous). Let E be the inductive limit of the system \((\mathrm{E}_{n},u_{n})\) (II, p.~29), and let \(v_{n}\) be the canonical mapping from \(E_{n}\) into E. Then the locally convex space E is Hausdorff. Moreover, for every subset A of E, the following conditions are equivalent :

\begin{enumerate}
\def\labelenumi{(\roman{enumi})}
\item
  A is bounded;
\item
  there exists an integer \(n\) such that \(\mathbf{A}\) is the image under \(v_{n}\) of a bounded subset of \(\mathrm{E}_n\) ;
\item
  A is relatively compact.
\end{enumerate}

We identify \(\mathbf{E}_n\) with a vector subspace of \(\mathbf{E}\) (endowed with a topology finer than the induced topology).

Lemma 1. --- Under the hypothesis of prop. 7, the topology of E is the finest topology for which all the mappings \(v_{n} : \mathrm{E}_{n} \to \mathrm{E}\) are continuous.

We need to prove that, if U is a subset of E such that \(U \cap E_{n}\) is open in \(E_{n}\) for every n, then U is open in E; in other words, we must prove that, for every \(x \in U\) , there exists a convex balanced set V such that \(x + V \subset U\) and that \(V \cap E_{n}\) is a neighbourhood of 0 in \(E_{n}\) for every large enough n (II, p.~27, prop. 5). For every n, let \(W_{n}\) be a convex balanced neighbourhood of 0 in \(E_{n}\) such that the closure \(H_{n}\) of \(W_{n}\) in \(E_{n+1}\) is compact. Let \(x \in U\) and let \(n_{0}\) be an integer such that \(x \in E_{n_{0}}\) . We shall construct, by induction, a sequence \((\varepsilon_{n})_{n \geqslant 0}\) of scalars \textgreater0 such that \(x + \sum_{n_{0} \leqslant i \leqslant n} \varepsilon_{i}H_{i}\) is contained in U for \(n \geqslant n_{0}\) . Suppose that the \(\varepsilon_{i}\) for i \textless{} n have been constructed. If \(n = n_{0}\) , set \(V_{n-1} = \{0\}\) ; if not, set

\[
\mathrm{V} _ {n - 1} = \sum_ {n _ {0} \leqslant i \leqslant n - 1} \varepsilon_ {i} \mathrm{H} _ {i}.
\]

Then \(V_{n-1}\) is compact in \(E_{n}\) , and a fortiori in \(E_{n+1}\) . Since \(U \cap E_{n+1}\) is open in \(E_{n+1}\) , there exists a scalar \(\varepsilon_{n} > 0\) such that \(x + V_{n} = x + V_{n-1} + \varepsilon_{n}H_{n}\) is contained in U (GT, II, §4, No.~3, cor.). Let \(V = \bigcup_{n \geqslant n_{0}} V_{n}\) . Then V is convex and balanced; for \(n \geqslant n_{0}\) , we have \(V \cap E_{n} \supset \varepsilon_{n}H_{n} \cap E_{n} \supset \varepsilon_{n}W_{n}\) , hence \(V \cap E_{n}\) is a neighbourhood of 0 in \(E_{n}\) . This completes the proof of the lemma.

The set \(\mathrm{U} = \mathrm{E} - \{0\}\) is such that the set \(\mathrm{U} \cap \mathrm{E}_n = \mathrm{E}_n - \{0\}\) is open in \(\mathrm{E}_n\) for every \(n\) , hence \(\mathrm{U}\) is open in \(\mathrm{E}\) , which proves that \(\mathrm{E}\) is Hausdorff (GT, III, § 1, No.~3, prop. 2). It is clear that property (ii) implies (iii) and that (iii) implies (i). Finally we show that (i) implies (ii). For this it is enough to show that if a subset \(\mathrm{A}\) of \(\mathrm{E}\) is not absorbed by any of the sets \(\sum_{0 \leqslant i \leqslant n} \mathrm{H}_i\) , then \(\mathrm{A}\) is not bounded. But then there exists a sequence \((x_n)_{n \geqslant 1}\) of points of \(\mathrm{A}\) such that, for every \(n\) , we have \(x_n \notin n^2 \sum_{0 \leqslant i \leqslant n} \mathrm{H}_i\) .

Then the set of the \(x_{n}/n\) is closed by virtue of lemma 1, since its intersection with \(E_{m}\) is discrete for every integer m. The complement of the set of the \(x_{n}/n\) is then an open neighbourhood of 0 which does not absorb the sequence \((x_{n})\) , hence A is not bounded.

Remarks. --- 1) With the above notations, let \(F_{n}\) be the vector space generated by \(H_{n}\) , with a norm equal to the gauge of \(H_{n}\) . We shall see (III, p.~8, cor.) that \(F_{n}\) is a Banach space. The injection from \(F_{n}\) into \(E_{n+1}\) is compact, hence a fortiori also the injection \(w_{n}\) from \(F_{n}\) into \(F_{n+1}\) . Further, E is the inductive limit of the inductive system ( \(F_{n}, w_{n}\) ) of Banach spaces. For, a convex balanced neighbourhood V of 0 in E is such that \(V \cap E_{n}\) absorbs \(H_{n-1}\) for all \(n \geqslant 1\) , and conversely, if a convex balanced set W in E is such that \(W \cap E_{n}\) absorbs \(H_{n-1}\) , then \(W \cap E_{n-1}\) contains a set homothetic to \(W_{n-1}\) for all \(n \geqslant 1\) , and hence W is a neighbourhood of 0 in E.

\begin{enumerate}
\def\labelenumi{\arabic{enumi})}
\setcounter{enumi}{1}
\tightlist
\item
  Let F be a locally convex space, k an integer \(\geqslant0\) and \(f:E^{k}\rightarrow F\) a multilinear mapping. For f to be continuous, it is necessary and sufficient that the restriction of f to \(E_{n}^{k}\) is continuous for every n.~We verify immediately that \(E^{k}\) has the final locally convex topology for the family of linear mappings \(v_{n}\times\cdots\times v_{n}:E_{n}\times\cdots\times E_{n}\to E\times\cdots\times E\) (II, p.~28, cor. 2 and p.~30, prop. 7) and that \(u_{n}\times\cdots\times u_{n}\) is a compact injective linear mapping from \((\mathrm{E}_{n})^{k}\) into \((\mathrm{E}_{n+1})^{k}\) . It is now enough to apply lemma 1.
\end{enumerate}
% semantic-fragments: legacy-input-seam
\relax{}
